%%%%%%%%%%%%%%%%%%%%%%%%%%%%%%%%%%%%%%%%%%%%%%%%%%%%%%%%%%%%
% 4.1 POINT-SET TOPOLOGY
% The Composition of Advanced Mathematics
%%%%%%%%%%%%%%%%%%%%%%%%%%%%%%%%%%%%%%%%%%%%%%%%%%%%%%%%%%%%

\section{Point-Set Topology}
\label{sec:point-set-topology}

\prerequisite{
Set Theory, Functions and Relations, Mathematical Logic,
Proof Techniques, Euclidean Geometry, and basic Analysis.
}

\learningobjectives{
By the end of this section, the reader should be able to:
\begin{itemize}
    \item explain what topology studies and how it differs from metric geometry;
    \item define a topological space using open sets;
    \item verify the topology axioms;
    \item identify trivial, discrete, indiscrete, cofinite, and metric topologies;
    \item compare finer and coarser topologies;
    \item define neighborhoods and neighborhood bases;
    \item determine interior, exterior, boundary, and closure of subsets;
    \item identify limit points, isolated points, and dense subsets;
    \item define closed sets and use their equivalent characterizations;
    \item explain bases and subbases for topologies;
    \item construct subspace, product, and quotient topologies;
    \item define continuity topologically;
    \item recognize homeomorphisms and topological invariants;
    \item distinguish convergence in topological and metric spaces;
    \item explain first and second countability;
    \item define compactness and recognize basic compactness results;
    \item define connectedness and path-connectedness;
    \item distinguish common separation axioms;
    \item explain Hausdorff spaces;
    \item recognize the role of point-set topology in analysis,
    geometry, manifolds, and later topology.
\end{itemize}
}

%%%%%%%%%%%%%%%%%%%%%%%%%%%%%%%%%%%%%%%%%%%%%%%%%%%%%%%%%%%%
\subsection{From Geometry to Topology}
\label{subsec:from-geometry-to-topology}
%%%%%%%%%%%%%%%%%%%%%%%%%%%%%%%%%%%%%%%%%%%%%%%%%%%%%%%%%%%%

Geometry studies properties such as

\[
\boxed{
\text{distance},
\quad
\text{angle},
\quad
\text{length},
\quad
\text{curvature}.
}
\]

Topology asks what remains when many of these metric measurements are
discarded.

The central ideas become

\[
\boxed{
\text{openness},
\quad
\text{continuity},
\quad
\text{connectedness},
\quad
\text{compactness}.
}
\]

\begin{conceptbox}
Topology studies properties of spaces that are preserved by continuous
deformations.

Distances and angles may change dramatically while the underlying
topological structure remains unchanged.
\end{conceptbox}

%%%%%%%%%%%%%%%%%%%%%%%%%%%%%%%%%%%%%%%%%%%%%%%%%%%%%%%%%%%%
\subsection{Topological Equivalence}
\label{subsec:topological-equivalence-intro}
%%%%%%%%%%%%%%%%%%%%%%%%%%%%%%%%%%%%%%%%%%%%%%%%%%%%%%%%%%%%

Two geometric objects can look very different while being topologically
equivalent.

For example, a circle can be stretched into an ellipse without cutting or
gluing.

From the viewpoint of Euclidean geometry, the two curves may have different
curvature.

From the viewpoint of topology, they have the same essential structure.

This motivates the study of \emph{homeomorphisms}, which will be defined
later in the section.

%%%%%%%%%%%%%%%%%%%%%%%%%%%%%%%%%%%%%%%%%%%%%%%%%%%%%%%%%%%%
\subsection{Open Sets as the Fundamental Objects}
\label{subsec:open-sets-fundamental}
%%%%%%%%%%%%%%%%%%%%%%%%%%%%%%%%%%%%%%%%%%%%%%%%%%%%%%%%%%%%

In metric spaces, open sets are usually introduced using open balls.

Topology reverses this perspective.

Instead of defining open sets from distance, we specify directly which
subsets are to be considered open.

This leads to the fundamental definition of a topological space.

%%%%%%%%%%%%%%%%%%%%%%%%%%%%%%%%%%%%%%%%%%%%%%%%%%%%%%%%%%%%
\subsection{Topological Spaces}
\label{subsec:topological-spaces}
%%%%%%%%%%%%%%%%%%%%%%%%%%%%%%%%%%%%%%%%%%%%%%%%%%%%%%%%%%%%

\begin{definitionbox}{Topology}
Let \(X\) be a set.

A \emph{topology} on \(X\) is a collection

\[
\boxed{
\mathcal{T}\subseteq\mathcal{P}(X)
}
\]

satisfying:

\begin{enumerate}
    \item
    \[
    \varnothing\in\mathcal{T}
    \qquad\text{and}\qquad
    X\in\mathcal{T};
    \]

    \item if
    \[
    \{U_\alpha\}_{\alpha\in A}
    \subseteq\mathcal{T},
    \]
    then
    \[
    \bigcup_{\alpha\in A}U_\alpha\in\mathcal{T};
    \]

    \item if
    \[
    U_1,\ldots,U_n\in\mathcal{T},
    \]
    then
    \[
    U_1\cap\cdots\cap U_n\in\mathcal{T}.
    \]
\end{enumerate}
\end{definitionbox}

The script letter

\[
\mathcal{T}
\]

is read ``script T'' and denotes the topology.

The symbol

\[
\mathcal{P}(X)
\]

denotes the power set of \(X\).

%%%%%%%%%%%%%%%%%%%%%%%%%%%%%%%%%%%%%%%%%%%%%%%%%%%%%%%%%%%%
\subsection{Topological Space}
\label{subsec:topological-space-pair}
%%%%%%%%%%%%%%%%%%%%%%%%%%%%%%%%%%%%%%%%%%%%%%%%%%%%%%%%%%%%

\begin{definitionbox}{Topological Space}
A \emph{topological space} is an ordered pair

\[
\boxed{
(X,\mathcal{T}),
}
\]

where \(X\) is a set and \(\mathcal{T}\) is a topology on \(X\).
\end{definitionbox}

Members of \(\mathcal{T}\) are called \emph{open sets}.

Thus

\[
\boxed{
U\subseteq X
\text{ is open}
\iff
U\in\mathcal{T}.
}
\]

%%%%%%%%%%%%%%%%%%%%%%%%%%%%%%%%%%%%%%%%%%%%%%%%%%%%%%%%%%%%
\subsection{Understanding the Topology Axioms}
\label{subsec:topology-axioms-meaning}
%%%%%%%%%%%%%%%%%%%%%%%%%%%%%%%%%%%%%%%%%%%%%%%%%%%%%%%%%%%%

The topology axioms say:

\[
\boxed{
\varnothing\text{ and }X\text{ are open};
}
\]

\[
\boxed{
\text{arbitrary unions of open sets are open};
}
\]

\[
\boxed{
\text{finite intersections of open sets are open}.
}
\]

The word \emph{arbitrary} is important.

A union may contain finitely many, countably many, or even uncountably many
open sets.

%%%%%%%%%%%%%%%%%%%%%%%%%%%%%%%%%%%%%%%%%%%%%%%%%%%%%%%%%%%%
\subsection{Why Only Finite Intersections?}
\label{subsec:finite-intersections-only}
%%%%%%%%%%%%%%%%%%%%%%%%%%%%%%%%%%%%%%%%%%%%%%%%%%%%%%%%%%%%

Infinite intersections of open sets need not be open.

For example, in the usual topology on \(\R\),

\[
U_n
=
\left(
-\frac1n,\frac1n
\right)
\]

is open for every positive integer \(n\).

However,

\[
\bigcap_{n=1}^{\infty}
\left(
-\frac1n,\frac1n
\right)
=
\{0\}.
\]

The singleton

\[
\{0\}
\]

is not open in the usual topology on \(\R\).

Therefore arbitrary intersections cannot be required to remain open.

%%%%%%%%%%%%%%%%%%%%%%%%%%%%%%%%%%%%%%%%%%%%%%%%%%%%%%%%%%%%
\subsection{The Indiscrete Topology}
\label{subsec:indiscrete-topology}
%%%%%%%%%%%%%%%%%%%%%%%%%%%%%%%%%%%%%%%%%%%%%%%%%%%%%%%%%%%%

\begin{definitionbox}{Indiscrete Topology}
The \emph{indiscrete topology} on a set \(X\) is

\[
\boxed{
\mathcal{T}
=
\{\varnothing,X\}.
}
\]
\end{definitionbox}

This is the smallest possible topology on \(X\).

Only the empty set and the entire space are open.

It is also called the \emph{trivial topology}.

%%%%%%%%%%%%%%%%%%%%%%%%%%%%%%%%%%%%%%%%%%%%%%%%%%%%%%%%%%%%
\subsection{The Discrete Topology}
\label{subsec:discrete-topology}
%%%%%%%%%%%%%%%%%%%%%%%%%%%%%%%%%%%%%%%%%%%%%%%%%%%%%%%%%%%%

\begin{definitionbox}{Discrete Topology}
The \emph{discrete topology} on \(X\) is

\[
\boxed{
\mathcal{T}
=
\mathcal{P}(X).
}
\]
\end{definitionbox}

Every subset of \(X\) is open.

This is the largest possible topology on \(X\).

%%%%%%%%%%%%%%%%%%%%%%%%%%%%%%%%%%%%%%%%%%%%%%%%%%%%%%%%%%%%
\subsection{Worked Example: A Finite Topology}
\label{subsec:worked-finite-topology}
%%%%%%%%%%%%%%%%%%%%%%%%%%%%%%%%%%%%%%%%%%%%%%%%%%%%%%%%%%%%

\begin{workedexamplebox}{Verify a Topology}

Let

\[
X=\{a,b,c\}
\]

and consider

\[
\mathcal{T}
=
\{
\varnothing,
\{a\},
\{a,b\},
X
\}.
\]

First,

\[
\varnothing,X\in\mathcal{T}.
\]

The possible unions of members of \(\mathcal{T}\) remain in
\(\mathcal{T}\).

For example,

\[
\{a\}\cup\{a,b\}
=
\{a,b\}.
\]

Finite intersections also remain in \(\mathcal{T}\):

\[
\{a\}\cap\{a,b\}
=
\{a\}.
\]

Therefore all topology axioms are satisfied.

\begin{answerbox}
\[
\boxed{
\mathcal{T}
\text{ is a topology on }X.
}
\]
\end{answerbox}
\end{workedexamplebox}

%%%%%%%%%%%%%%%%%%%%%%%%%%%%%%%%%%%%%%%%%%%%%%%%%%%%%%%%%%%%
\subsection{A Collection That Is Not a Topology}
\label{subsec:not-a-topology-example}
%%%%%%%%%%%%%%%%%%%%%%%%%%%%%%%%%%%%%%%%%%%%%%%%%%%%%%%%%%%%

Let

\[
X=\{a,b,c\}
\]

and consider

\[
\mathcal{S}
=
\{
\varnothing,
X,
\{a\},
\{b\}
\}.
\]

Although

\[
\{a\},\{b\}\in\mathcal{S},
\]

their union is

\[
\{a\}\cup\{b\}
=
\{a,b\}.
\]

But

\[
\{a,b\}\notin\mathcal{S}.
\]

Therefore

\[
\boxed{
\mathcal{S}
\text{ is not a topology}.
}
\]

%%%%%%%%%%%%%%%%%%%%%%%%%%%%%%%%%%%%%%%%%%%%%%%%%%%%%%%%%%%%
\subsection{Metric Topologies}
\label{subsec:metric-topologies}
%%%%%%%%%%%%%%%%%%%%%%%%%%%%%%%%%%%%%%%%%%%%%%%%%%%%%%%%%%%%

Let

\[
(X,d)
\]

be a metric space.

For \(x\in X\) and \(r>0\), the open ball centered at \(x\) with radius
\(r\) is

\[
\boxed{
B_r(x)
=
\{y\in X:d(x,y)<r\}.
}
\]

A subset \(U\subseteq X\) is declared open when every point \(x\in U\)
lies inside some open ball contained in \(U\).

That is,

\[
\boxed{
U\text{ open}
\iff
\forall x\in U,\ 
\exists r>0
\text{ such that }
B_r(x)\subseteq U.
}
\]

The resulting topology is called the \emph{metric topology induced by \(d\)}.

%%%%%%%%%%%%%%%%%%%%%%%%%%%%%%%%%%%%%%%%%%%%%%%%%%%%%%%%%%%%
\subsection{The Usual Topology on the Real Line}
\label{subsec:usual-topology-real-line}
%%%%%%%%%%%%%%%%%%%%%%%%%%%%%%%%%%%%%%%%%%%%%%%%%%%%%%%%%%%%

The standard metric

\[
d(x,y)=|x-y|
\]

on \(\R\) induces the usual topology.

Open intervals

\[
(a,b)
\]

are open because every point \(x\in(a,b)\) has a sufficiently small interval

\[
(x-r,x+r)
\]

contained in \((a,b)\).

The usual topology on \(\R\) can therefore be generated from open intervals.

%%%%%%%%%%%%%%%%%%%%%%%%%%%%%%%%%%%%%%%%%%%%%%%%%%%%%%%%%%%%
\subsection{The Euclidean Topology on \texorpdfstring{\(\R^n\)}{R^n}}
\label{subsec:euclidean-topology-rn}
%%%%%%%%%%%%%%%%%%%%%%%%%%%%%%%%%%%%%%%%%%%%%%%%%%%%%%%%%%%%

The Euclidean metric on \(\R^n\) is

\[
\boxed{
d(x,y)
=
\|x-y\|
=
\sqrt{
\sum_{i=1}^{n}(x_i-y_i)^2
}.
}
\]

Its open balls generate the standard Euclidean topology on \(\R^n\).

Thus ordinary multivariable calculus takes place inside a particular
topological space.

%%%%%%%%%%%%%%%%%%%%%%%%%%%%%%%%%%%%%%%%%%%%%%%%%%%%%%%%%%%%
\subsection{Different Metrics Can Produce the Same Topology}
\label{subsec:different-metrics-same-topology}
%%%%%%%%%%%%%%%%%%%%%%%%%%%%%%%%%%%%%%%%%%%%%%%%%%%%%%%%%%%%

The metric determines distances, while the induced topology remembers only
which sets are open.

Different metrics can therefore produce the same topology.

For example, on \(\R^n\), the Euclidean metric

\[
d_2(x,y)
=
\left(
\sum_{i=1}^{n}|x_i-y_i|^2
\right)^{1/2}
\]

and the taxicab metric

\[
d_1(x,y)
=
\sum_{i=1}^{n}|x_i-y_i|
\]

generate the same standard topology.

Their metric balls have different shapes, but they determine the same
notion of openness.

%%%%%%%%%%%%%%%%%%%%%%%%%%%%%%%%%%%%%%%%%%%%%%%%%%%%%%%%%%%%
\subsection{Topological Versus Metric Information}
\label{subsec:topological-vs-metric-information}
%%%%%%%%%%%%%%%%%%%%%%%%%%%%%%%%%%%%%%%%%%%%%%%%%%%%%%%%%%%%

A topology does not generally tell us the numerical distance between two
points.

It tells us which neighborhoods surround each point.

Thus topology retains enough information to study continuity while
discarding much of the metric structure.

\[
\boxed{
\text{metric}
\Longrightarrow
\text{topology}
}
\]

but not every topology necessarily arises from a metric.

%%%%%%%%%%%%%%%%%%%%%%%%%%%%%%%%%%%%%%%%%%%%%%%%%%%%%%%%%%%%
\subsection{Metrizable Spaces}
\label{subsec:metrizable-spaces}
%%%%%%%%%%%%%%%%%%%%%%%%%%%%%%%%%%%%%%%%%%%%%%%%%%%%%%%%%%%%

\begin{definitionbox}{Metrizable Space}
A topological space is \emph{metrizable} if there exists a metric whose
induced topology is the given topology.
\end{definitionbox}

Determining when an abstract topology is metrizable is an important question
in point-set topology.

%%%%%%%%%%%%%%%%%%%%%%%%%%%%%%%%%%%%%%%%%%%%%%%%%%%%%%%%%%%%
\subsection{The Cofinite Topology}
\label{subsec:cofinite-topology}
%%%%%%%%%%%%%%%%%%%%%%%%%%%%%%%%%%%%%%%%%%%%%%%%%%%%%%%%%%%%

Let \(X\) be a set.

The \emph{cofinite topology} is

\[
\boxed{
\mathcal{T}_{\mathrm{cof}}
=
\{\varnothing\}
\cup
\{
U\subseteq X:X\setminus U\text{ is finite}
\}.
}
\]

Thus every nonempty open set has finite complement.

If \(X\) is infinite, this gives a topology quite different from the
discrete topology.

%%%%%%%%%%%%%%%%%%%%%%%%%%%%%%%%%%%%%%%%%%%%%%%%%%%%%%%%%%%%
\subsection{Finer and Coarser Topologies}
\label{subsec:finer-coarser-topologies}
%%%%%%%%%%%%%%%%%%%%%%%%%%%%%%%%%%%%%%%%%%%%%%%%%%%%%%%%%%%%

Suppose

\[
\mathcal{T}_1
\]

and

\[
\mathcal{T}_2
\]

are topologies on the same set \(X\).

\begin{definitionbox}{Finer Topology}
If

\[
\boxed{
\mathcal{T}_1\subseteq\mathcal{T}_2,
}
\]

then \(\mathcal{T}_2\) is called \emph{finer} than \(\mathcal{T}_1\).

Equivalently, \(\mathcal{T}_1\) is \emph{coarser} than \(\mathcal{T}_2\).
\end{definitionbox}

A finer topology contains more open sets.

%%%%%%%%%%%%%%%%%%%%%%%%%%%%%%%%%%%%%%%%%%%%%%%%%%%%%%%%%%%%
\subsection{Extreme Topologies}
\label{subsec:extreme-topologies}
%%%%%%%%%%%%%%%%%%%%%%%%%%%%%%%%%%%%%%%%%%%%%%%%%%%%%%%%%%%%

For every topology \(\mathcal{T}\) on \(X\),

\[
\boxed{
\{\varnothing,X\}
\subseteq
\mathcal{T}
\subseteq
\mathcal{P}(X).
}
\]

Therefore

\[
\boxed{
\text{indiscrete topology}
\leq
\mathcal{T}
\leq
\text{discrete topology},
}
\]

where the ordering is by inclusion.

%%%%%%%%%%%%%%%%%%%%%%%%%%%%%%%%%%%%%%%%%%%%%%%%%%%%%%%%%%%%
\subsection{Closed Sets}
\label{subsec:closed-sets}
%%%%%%%%%%%%%%%%%%%%%%%%%%%%%%%%%%%%%%%%%%%%%%%%%%%%%%%%%%%%

\begin{definitionbox}{Closed Set}
A subset \(F\subseteq X\) is \emph{closed} if its complement

\[
\boxed{
X\setminus F
}
\]

is open.
\end{definitionbox}

Thus

\[
\boxed{
F\text{ closed}
\iff
X\setminus F\in\mathcal{T}.
}
\]

%%%%%%%%%%%%%%%%%%%%%%%%%%%%%%%%%%%%%%%%%%%%%%%%%%%%%%%%%%%%
\subsection{Closed-Set Axioms}
\label{subsec:closed-set-axioms}
%%%%%%%%%%%%%%%%%%%%%%%%%%%%%%%%%%%%%%%%%%%%%%%%%%%%%%%%%%%%

By taking complements of the open-set axioms, we obtain:

\[
\boxed{
\varnothing\text{ and }X\text{ are closed};
}
\]

\[
\boxed{
\text{arbitrary intersections of closed sets are closed};
}
\]

\[
\boxed{
\text{finite unions of closed sets are closed}.
}
\]

These statements follow from De Morgan's laws.

%%%%%%%%%%%%%%%%%%%%%%%%%%%%%%%%%%%%%%%%%%%%%%%%%%%%%%%%%%%%
\subsection{Sets Can Be Both Open and Closed}
\label{subsec:clopen-sets}
%%%%%%%%%%%%%%%%%%%%%%%%%%%%%%%%%%%%%%%%%%%%%%%%%%%%%%%%%%%%

A set may be both open and closed.

Such a set is called \emph{clopen}.

For every topological space,

\[
\boxed{
\varnothing
\quad\text{and}\quad
X
}
\]

are clopen.

In a discrete space, every subset is clopen.

\begin{warningbox}
``Closed'' does not mean ``not open.''

A set may be open, closed, both, or neither.
\end{warningbox}

%%%%%%%%%%%%%%%%%%%%%%%%%%%%%%%%%%%%%%%%%%%%%%%%%%%%%%%%%%%%
\subsection{Neighborhoods}
\label{subsec:topological-neighborhoods}
%%%%%%%%%%%%%%%%%%%%%%%%%%%%%%%%%%%%%%%%%%%%%%%%%%%%%%%%%%%%

\begin{definitionbox}{Neighborhood}
A set \(N\subseteq X\) is a \emph{neighborhood} of \(x\in X\) if there exists
an open set \(U\) such that

\[
\boxed{
x\in U\subseteq N.
}
\]
\end{definitionbox}

A neighborhood need not itself be open.

An \emph{open neighborhood} is simply an open set containing the point.

%%%%%%%%%%%%%%%%%%%%%%%%%%%%%%%%%%%%%%%%%%%%%%%%%%%%%%%%%%%%
\subsection{Neighborhood Systems}
\label{subsec:neighborhood-systems}
%%%%%%%%%%%%%%%%%%%%%%%%%%%%%%%%%%%%%%%%%%%%%%%%%%%%%%%%%%%%

The collection of all neighborhoods of \(x\) is often denoted

\[
\boxed{
\mathcal{N}(x).
}
\]

Knowledge of all neighborhood systems determines the topology.

Thus topology can be described either through open sets or through the local
neighborhood structure surrounding each point.

%%%%%%%%%%%%%%%%%%%%%%%%%%%%%%%%%%%%%%%%%%%%%%%%%%%%%%%%%%%%
\subsection{Interior}
\label{subsec:topological-interior}
%%%%%%%%%%%%%%%%%%%%%%%%%%%%%%%%%%%%%%%%%%%%%%%%%%%%%%%%%%%%

\begin{definitionbox}{Interior}
Let \(A\subseteq X\).

The \emph{interior} of \(A\), denoted

\[
\boxed{
\operatorname{int}(A)
}
\]

or \(A^\circ\), is the union of all open subsets of \(A\).
\end{definitionbox}

Equivalently,

\[
\boxed{
x\in\operatorname{int}(A)
\iff
\exists U\text{ open such that }
x\in U\subseteq A.
}
\]

%%%%%%%%%%%%%%%%%%%%%%%%%%%%%%%%%%%%%%%%%%%%%%%%%%%%%%%%%%%%
\subsection{Properties of Interior}
\label{subsec:interior-properties}
%%%%%%%%%%%%%%%%%%%%%%%%%%%%%%%%%%%%%%%%%%%%%%%%%%%%%%%%%%%%

For every \(A\subseteq X\),

\[
\boxed{
\operatorname{int}(A)\subseteq A.
}
\]

Furthermore,

\[
\boxed{
A\text{ is open}
\iff
\operatorname{int}(A)=A.
}
\]

The interior is the largest open subset contained in \(A\).

%%%%%%%%%%%%%%%%%%%%%%%%%%%%%%%%%%%%%%%%%%%%%%%%%%%%%%%%%%%%
\subsection{Closure}
\label{subsec:topological-closure}
%%%%%%%%%%%%%%%%%%%%%%%%%%%%%%%%%%%%%%%%%%%%%%%%%%%%%%%%%%%%

\begin{definitionbox}{Closure}
The \emph{closure} of \(A\subseteq X\), denoted

\[
\boxed{
\overline{A},
}
\]

is the intersection of all closed subsets of \(X\) containing \(A\).
\end{definitionbox}

Therefore \(\overline{A}\) is the smallest closed set containing \(A\).

In particular,

\[
\boxed{
A\subseteq\overline{A}.
}
\]

%%%%%%%%%%%%%%%%%%%%%%%%%%%%%%%%%%%%%%%%%%%%%%%%%%%%%%%%%%%%
\subsection{Neighborhood Characterization of Closure}
\label{subsec:closure-neighborhood-characterization}
%%%%%%%%%%%%%%%%%%%%%%%%%%%%%%%%%%%%%%%%%%%%%%%%%%%%%%%%%%%%

A point \(x\) belongs to \(\overline{A}\) precisely when every open
neighborhood of \(x\) intersects \(A\):

\[
\boxed{
x\in\overline{A}
\iff
\forall U\ni x,\ 
U\text{ open}
\Longrightarrow
U\cap A\neq\varnothing.
}
\]

This is one of the most useful characterizations of closure.

%%%%%%%%%%%%%%%%%%%%%%%%%%%%%%%%%%%%%%%%%%%%%%%%%%%%%%%%%%%%
\subsection{Properties of Closure}
\label{subsec:closure-properties}
%%%%%%%%%%%%%%%%%%%%%%%%%%%%%%%%%%%%%%%%%%%%%%%%%%%%%%%%%%%%

For subsets \(A,B\subseteq X\),

\[
\boxed{
A\subseteq\overline{A},
}
\]

\[
\boxed{
\overline{\overline{A}}
=
\overline{A},
}
\]

and

\[
\boxed{
A\subseteq B
\Longrightarrow
\overline{A}\subseteq\overline{B}.
}
\]

Also,

\[
\boxed{
A\text{ closed}
\iff
\overline{A}=A.
}
\]

%%%%%%%%%%%%%%%%%%%%%%%%%%%%%%%%%%%%%%%%%%%%%%%%%%%%%%%%%%%%
\subsection{Boundary}
\label{subsec:topological-boundary}
%%%%%%%%%%%%%%%%%%%%%%%%%%%%%%%%%%%%%%%%%%%%%%%%%%%%%%%%%%%%

\begin{definitionbox}{Boundary}
The \emph{boundary} of \(A\subseteq X\) is

\[
\boxed{
\partial A
=
\overline{A}
\setminus
\operatorname{int}(A).
}
\]
\end{definitionbox}

The symbol

\[
\partial
\]

is read ``partial'' or, in this context, ``boundary.''

A point lies on the boundary of \(A\) when every neighborhood of the point
meets both \(A\) and its complement.

%%%%%%%%%%%%%%%%%%%%%%%%%%%%%%%%%%%%%%%%%%%%%%%%%%%%%%%%%%%%
\subsection{Exterior}
\label{subsec:topological-exterior}
%%%%%%%%%%%%%%%%%%%%%%%%%%%%%%%%%%%%%%%%%%%%%%%%%%%%%%%%%%%%

The exterior of \(A\) is

\[
\boxed{
\operatorname{ext}(A)
=
\operatorname{int}(X\setminus A).
}
\]

The space decomposes into three disjoint parts:

\[
\boxed{
X
=
\operatorname{int}(A)
\cup
\partial A
\cup
\operatorname{ext}(A).
}
\]

%%%%%%%%%%%%%%%%%%%%%%%%%%%%%%%%%%%%%%%%%%%%%%%%%%%%%%%%%%%%
\subsection{Worked Example: Interior, Closure, and Boundary}
\label{subsec:worked-interior-closure-boundary}
%%%%%%%%%%%%%%%%%%%%%%%%%%%%%%%%%%%%%%%%%%%%%%%%%%%%%%%%%%%%

\begin{workedexamplebox}{An Interval in the Real Line}

Let

\[
A=(0,1]
\subseteq\R
\]

with the usual topology.

The interior is

\[
\operatorname{int}(A)
=
(0,1).
\]

The closure is

\[
\overline{A}
=
[0,1].
\]

Therefore the boundary is

\[
\partial A
=
[0,1]\setminus(0,1)
=
\{0,1\}.
\]

The exterior is

\[
(-\infty,0)\cup(1,\infty).
\]

\begin{answerbox}
\[
\boxed{
\operatorname{int}(A)=(0,1),
\qquad
\overline{A}=[0,1],
\qquad
\partial A=\{0,1\}.
}
\]
\end{answerbox}
\end{workedexamplebox}

%%%%%%%%%%%%%%%%%%%%%%%%%%%%%%%%%%%%%%%%%%%%%%%%%%%%%%%%%%%%
\subsection{Limit Points}
\label{subsec:topological-limit-points}
%%%%%%%%%%%%%%%%%%%%%%%%%%%%%%%%%%%%%%%%%%%%%%%%%%%%%%%%%%%%

\begin{definitionbox}{Limit Point}
A point \(x\in X\) is a \emph{limit point}, or \emph{accumulation point}, of
\(A\subseteq X\) if every neighborhood of \(x\) contains a point of
\(A\) different from \(x\).
\end{definitionbox}

Equivalently,

\[
\boxed{
U\cap(A\setminus\{x\})\neq\varnothing
}
\]

for every open neighborhood \(U\) of \(x\).

%%%%%%%%%%%%%%%%%%%%%%%%%%%%%%%%%%%%%%%%%%%%%%%%%%%%%%%%%%%%
\subsection{Derived Set}
\label{subsec:derived-set}
%%%%%%%%%%%%%%%%%%%%%%%%%%%%%%%%%%%%%%%%%%%%%%%%%%%%%%%%%%%%

The set of all limit points of \(A\) is called the \emph{derived set} and is
often denoted

\[
\boxed{
A'.
}
\]

In many familiar spaces,

\[
\boxed{
\overline{A}=A\cup A'.
}
\]

More generally, this identity holds in every topological space under the
standard neighborhood definition of limit point.

%%%%%%%%%%%%%%%%%%%%%%%%%%%%%%%%%%%%%%%%%%%%%%%%%%%%%%%%%%%%
\subsection{Isolated Points}
\label{subsec:isolated-points}
%%%%%%%%%%%%%%%%%%%%%%%%%%%%%%%%%%%%%%%%%%%%%%%%%%%%%%%%%%%%

\begin{definitionbox}{Isolated Point}
A point \(x\in A\) is \emph{isolated in \(A\)} if there exists an open set
\(U\) such that

\[
\boxed{
U\cap A=\{x\}.
}
\]
\end{definitionbox}

Thus an isolated point belongs to the set but is not a limit point of the
set.

%%%%%%%%%%%%%%%%%%%%%%%%%%%%%%%%%%%%%%%%%%%%%%%%%%%%%%%%%%%%
\subsection{Example: The Integers in the Real Line}
\label{subsec:integers-isolated}
%%%%%%%%%%%%%%%%%%%%%%%%%%%%%%%%%%%%%%%%%%%%%%%%%%%%%%%%%%%%

Consider

\[
\Z\subseteq\R
\]

with the usual topology.

Every integer \(n\) is isolated in \(\Z\), because

\[
\left(
n-\frac12,
n+\frac12
\right)
\cap\Z
=
\{n\}.
\]

Therefore \(\Z\) has no limit points in \(\R\).

It is consequently closed in \(\R\).

%%%%%%%%%%%%%%%%%%%%%%%%%%%%%%%%%%%%%%%%%%%%%%%%%%%%%%%%%%%%
\subsection{Dense Sets}
\label{subsec:dense-sets}
%%%%%%%%%%%%%%%%%%%%%%%%%%%%%%%%%%%%%%%%%%%%%%%%%%%%%%%%%%%%

\begin{definitionbox}{Dense Subset}
A subset \(A\subseteq X\) is \emph{dense in \(X\)} if

\[
\boxed{
\overline{A}=X.
}
\]
\end{definitionbox}

Equivalently, every nonempty open subset of \(X\) intersects \(A\).

%%%%%%%%%%%%%%%%%%%%%%%%%%%%%%%%%%%%%%%%%%%%%%%%%%%%%%%%%%%%
\subsection{The Rational Numbers Are Dense}
\label{subsec:rationals-dense}
%%%%%%%%%%%%%%%%%%%%%%%%%%%%%%%%%%%%%%%%%%%%%%%%%%%%%%%%%%%%

In the usual topology on \(\R\),

\[
\boxed{
\overline{\Q}=\R.
}
\]

Every open interval contains a rational number.

Similarly, the irrational numbers are dense in \(\R\).

Thus two disjoint subsets can both be dense in the same space.

%%%%%%%%%%%%%%%%%%%%%%%%%%%%%%%%%%%%%%%%%%%%%%%%%%%%%%%%%%%%
\subsection{Bases for Topologies}
\label{subsec:topological-bases}
%%%%%%%%%%%%%%%%%%%%%%%%%%%%%%%%%%%%%%%%%%%%%%%%%%%%%%%%%%%%

Listing every open set can be inconvenient.

A smaller collection can often generate the topology.

\begin{definitionbox}{Basis}
A collection \(\mathcal{B}\) of subsets of \(X\) is a \emph{basis} for a
topology if:

\begin{enumerate}
    \item for every \(x\in X\), some \(B\in\mathcal{B}\) contains \(x\);

    \item whenever
    \[
    x\in B_1\cap B_2
    \]
    with \(B_1,B_2\in\mathcal{B}\), there exists \(B_3\in\mathcal{B}\)
    such that
    \[
    x\in B_3\subseteq B_1\cap B_2.
    \]
\end{enumerate}
\end{definitionbox}

%%%%%%%%%%%%%%%%%%%%%%%%%%%%%%%%%%%%%%%%%%%%%%%%%%%%%%%%%%%%
\subsection{Topology Generated by a Basis}
\label{subsec:topology-generated-by-basis}
%%%%%%%%%%%%%%%%%%%%%%%%%%%%%%%%%%%%%%%%%%%%%%%%%%%%%%%%%%%%

The topology generated by \(\mathcal{B}\) consists of arbitrary unions of
basis elements.

Thus

\[
\boxed{
U\text{ open}
\iff
U=\bigcup_{\alpha\in A}B_\alpha
}
\]

for suitable

\[
B_\alpha\in\mathcal{B}.
\]

Equivalently,

\[
\boxed{
U\text{ open}
\iff
\forall x\in U,\
\exists B\in\mathcal{B}
\text{ with }
x\in B\subseteq U.
}
\]

%%%%%%%%%%%%%%%%%%%%%%%%%%%%%%%%%%%%%%%%%%%%%%%%%%%%%%%%%%%%
\subsection{Standard Basis for the Real Line}
\label{subsec:standard-basis-real-line}
%%%%%%%%%%%%%%%%%%%%%%%%%%%%%%%%%%%%%%%%%%%%%%%%%%%%%%%%%%%%

The collection

\[
\boxed{
\mathcal{B}
=
\{(a,b):a<b\}
}
\]

forms a basis for the usual topology on \(\R\).

In fact, it is enough to use intervals with rational endpoints:

\[
\boxed{
\mathcal{B}_{\Q}
=
\{(a,b):a,b\in\Q,\ a<b\}.
}
\]

This basis is countable.

%%%%%%%%%%%%%%%%%%%%%%%%%%%%%%%%%%%%%%%%%%%%%%%%%%%%%%%%%%%%
\subsection{Subbases}
\label{subsec:topological-subbases}
%%%%%%%%%%%%%%%%%%%%%%%%%%%%%%%%%%%%%%%%%%%%%%%%%%%%%%%%%%%%

\begin{definitionbox}{Subbasis}
A collection \(\mathcal{S}\) of subsets of \(X\) is a \emph{subbasis} if
finite intersections of elements of \(\mathcal{S}\) form a basis for a
topology on \(X\).
\end{definitionbox}

The topology generated by \(\mathcal{S}\) consists of arbitrary unions of
finite intersections of subbasis elements.

%%%%%%%%%%%%%%%%%%%%%%%%%%%%%%%%%%%%%%%%%%%%%%%%%%%%%%%%%%%%
\subsection{Local Bases}
\label{subsec:local-bases}
%%%%%%%%%%%%%%%%%%%%%%%%%%%%%%%%%%%%%%%%%%%%%%%%%%%%%%%%%%%%

A \emph{local basis} at a point \(x\) is a collection of neighborhoods of
\(x\) such that every neighborhood of \(x\) contains one of them.

In a metric space,

\[
\boxed{
\left\{
B_{1/n}(x):
n\in\mathbb{N}
\right\}
}
\]

forms a countable local basis at \(x\).

%%%%%%%%%%%%%%%%%%%%%%%%%%%%%%%%%%%%%%%%%%%%%%%%%%%%%%%%%%%%
\subsection{First Countability}
\label{subsec:first-countability}
%%%%%%%%%%%%%%%%%%%%%%%%%%%%%%%%%%%%%%%%%%%%%%%%%%%%%%%%%%%%

\begin{definitionbox}{First Countable}
A topological space is \emph{first countable} if every point has a countable
local basis.
\end{definitionbox}

Every metric space is first countable.

First countability helps sequences describe local topological behavior.

%%%%%%%%%%%%%%%%%%%%%%%%%%%%%%%%%%%%%%%%%%%%%%%%%%%%%%%%%%%%
\subsection{Second Countability}
\label{subsec:second-countability}
%%%%%%%%%%%%%%%%%%%%%%%%%%%%%%%%%%%%%%%%%%%%%%%%%%%%%%%%%%%%

\begin{definitionbox}{Second Countable}
A topological space is \emph{second countable} if its topology has a
countable basis.
\end{definitionbox}

The usual topology on \(\R\) is second countable because intervals with
rational endpoints form a countable basis.

Similarly,

\[
\R^n
\]

is second countable.

%%%%%%%%%%%%%%%%%%%%%%%%%%%%%%%%%%%%%%%%%%%%%%%%%%%%%%%%%%%%
\subsection{Subspace Topology}
\label{subsec:subspace-topology}
%%%%%%%%%%%%%%%%%%%%%%%%%%%%%%%%%%%%%%%%%%%%%%%%%%%%%%%%%%%%

Let

\[
Y\subseteq X,
\]

where \(X\) is a topological space.

\begin{definitionbox}{Subspace Topology}
The \emph{subspace topology} on \(Y\) is

\[
\boxed{
\mathcal{T}_Y
=
\{
Y\cap U:
U\in\mathcal{T}
\}.
}
\]
\end{definitionbox}

Thus a subset \(V\subseteq Y\) is open in \(Y\) exactly when

\[
\boxed{
V=Y\cap U
}
\]

for some open set \(U\subseteq X\).

%%%%%%%%%%%%%%%%%%%%%%%%%%%%%%%%%%%%%%%%%%%%%%%%%%%%%%%%%%%%
\subsection{Relative Openness}
\label{subsec:relative-openness}
%%%%%%%%%%%%%%%%%%%%%%%%%%%%%%%%%%%%%%%%%%%%%%%%%%%%%%%%%%%%

A set can be open in a subspace without being open in the ambient space.

For example, let

\[
Y=[0,1]\subseteq\R.
\]

Then

\[
[0,1/2)
\]

is open in \(Y\), because

\[
\boxed{
[0,1/2)
=
Y\cap(-1,1/2).
}
\]

However,

\[
[0,1/2)
\]

is not open in \(\R\).

%%%%%%%%%%%%%%%%%%%%%%%%%%%%%%%%%%%%%%%%%%%%%%%%%%%%%%%%%%%%
\subsection{Product Topology}
\label{subsec:product-topology}
%%%%%%%%%%%%%%%%%%%%%%%%%%%%%%%%%%%%%%%%%%%%%%%%%%%%%%%%%%%%

Let \(X\) and \(Y\) be topological spaces.

The product set is

\[
X\times Y.
\]

The product topology has basis

\[
\boxed{
\mathcal{B}
=
\{
U\times V:
U\subseteq X\text{ open},
\;
V\subseteq Y\text{ open}
\}.
}
\]

Thus open rectangles generate the product topology.

%%%%%%%%%%%%%%%%%%%%%%%%%%%%%%%%%%%%%%%%%%%%%%%%%%%%%%%%%%%%
\subsection{Example: Product Topology on the Plane}
\label{subsec:product-topology-plane}
%%%%%%%%%%%%%%%%%%%%%%%%%%%%%%%%%%%%%%%%%%%%%%%%%%%%%%%%%%%%

Since

\[
\R^2
=
\R\times\R,
\]

products of open intervals

\[
(a,b)\times(c,d)
\]

form a basis for the usual topology on \(\R^2\).

Although these basis elements are rectangles rather than Euclidean balls,
they generate the same topology.

%%%%%%%%%%%%%%%%%%%%%%%%%%%%%%%%%%%%%%%%%%%%%%%%%%%%%%%%%%%%
\subsection{Finite and Infinite Products}
\label{subsec:finite-infinite-products}
%%%%%%%%%%%%%%%%%%%%%%%%%%%%%%%%%%%%%%%%%%%%%%%%%%%%%%%%%%%%

For finitely many spaces,

\[
X_1\times\cdots\times X_n,
\]

basis elements are products

\[
U_1\times\cdots\times U_n
\]

with each \(U_i\) open.

For infinitely many spaces, the product topology requires that only finitely
many coordinates be restricted away from the entire coordinate space in a
basic open set.

This distinction becomes important in functional analysis and general
topology.

%%%%%%%%%%%%%%%%%%%%%%%%%%%%%%%%%%%%%%%%%%%%%%%%%%%%%%%%%%%%
\subsection{Quotient Topology}
\label{subsec:quotient-topology}
%%%%%%%%%%%%%%%%%%%%%%%%%%%%%%%%%%%%%%%%%%%%%%%%%%%%%%%%%%%%

Topology can also be constructed by identifying points.

Let

\[
q:X\to Y
\]

be a surjective map.

\begin{definitionbox}{Quotient Topology}
The \emph{quotient topology} on \(Y\) induced by \(q\) is defined by

\[
\boxed{
U\subseteq Y\text{ open}
\iff
q^{-1}(U)\text{ open in }X.
}
\]
\end{definitionbox}

The map \(q\) is called a \emph{quotient map} when \(Y\) carries this
topology.

%%%%%%%%%%%%%%%%%%%%%%%%%%%%%%%%%%%%%%%%%%%%%%%%%%%%%%%%%%%%
\subsection{Gluing Endpoints of an Interval}
\label{subsec:interval-to-circle-quotient}
%%%%%%%%%%%%%%%%%%%%%%%%%%%%%%%%%%%%%%%%%%%%%%%%%%%%%%%%%%%%

Consider the interval

\[
[0,1].
\]

Identify the endpoints

\[
0\sim1.
\]

The resulting quotient space is topologically a circle:

\[
\boxed{
[0,1]/(0\sim1)
\cong
S^1.
}
\]

Here

\[
\sim
\]

is read ``is equivalent to,'' and denotes an equivalence relation.

The symbol

\[
\cong
\]

is read ``is isomorphic to'' or, in this topological context,
``is homeomorphic to.''

%%%%%%%%%%%%%%%%%%%%%%%%%%%%%%%%%%%%%%%%%%%%%%%%%%%%%%%%%%%%
\subsection{Continuity}
\label{subsec:topological-continuity}
%%%%%%%%%%%%%%%%%%%%%%%%%%%%%%%%%%%%%%%%%%%%%%%%%%%%%%%%%%%%

The familiar epsilon-delta definition of continuity depends on distance.

Topology gives a more general definition.

\begin{definitionbox}{Continuous Map}
Let

\[
f:X\to Y
\]

be a function between topological spaces.

The function \(f\) is \emph{continuous} if

\[
\boxed{
f^{-1}(U)
\text{ is open in }X
}
\]

for every open set

\[
U\subseteq Y.
\]
\end{definitionbox}

%%%%%%%%%%%%%%%%%%%%%%%%%%%%%%%%%%%%%%%%%%%%%%%%%%%%%%%%%%%%
\subsection{Why Preimages Appear}
\label{subsec:continuity-preimages}
%%%%%%%%%%%%%%%%%%%%%%%%%%%%%%%%%%%%%%%%%%%%%%%%%%%%%%%%%%%%

Continuity is defined using preimages rather than images because preimages
preserve the set operations needed by topology:

\[
f^{-1}
\left(
\bigcup_\alpha U_\alpha
\right)
=
\bigcup_\alpha f^{-1}(U_\alpha),
\]

and

\[
f^{-1}
\left(
\bigcap_\alpha U_\alpha
\right)
=
\bigcap_\alpha f^{-1}(U_\alpha).
\]

Thus inverse images interact naturally with open-set axioms.

%%%%%%%%%%%%%%%%%%%%%%%%%%%%%%%%%%%%%%%%%%%%%%%%%%%%%%%%%%%%
\subsection{Equivalent Closed-Set Definition}
\label{subsec:continuity-closed-sets}
%%%%%%%%%%%%%%%%%%%%%%%%%%%%%%%%%%%%%%%%%%%%%%%%%%%%%%%%%%%%

A map

\[
f:X\to Y
\]

is continuous if and only if the preimage of every closed subset of \(Y\)
is closed in \(X\).

Therefore

\[
\boxed{
F\subseteq Y\text{ closed}
\Longrightarrow
f^{-1}(F)\subseteq X\text{ closed}.
}
\]

%%%%%%%%%%%%%%%%%%%%%%%%%%%%%%%%%%%%%%%%%%%%%%%%%%%%%%%%%%%%
\subsection{Composition of Continuous Maps}
\label{subsec:composition-continuous}
%%%%%%%%%%%%%%%%%%%%%%%%%%%%%%%%%%%%%%%%%%%%%%%%%%%%%%%%%%%%

If

\[
f:X\to Y
\]

and

\[
g:Y\to Z
\]

are continuous, then

\[
\boxed{
g\circ f:X\to Z
}
\]

is continuous.

Indeed,

\[
(g\circ f)^{-1}(U)
=
f^{-1}(g^{-1}(U)).
\]

Since inverse images of open sets remain open at each stage, the composition
is continuous.

%%%%%%%%%%%%%%%%%%%%%%%%%%%%%%%%%%%%%%%%%%%%%%%%%%%%%%%%%%%%
\subsection{Homeomorphisms}
\label{subsec:homeomorphisms}
%%%%%%%%%%%%%%%%%%%%%%%%%%%%%%%%%%%%%%%%%%%%%%%%%%%%%%%%%%%%

\begin{definitionbox}{Homeomorphism}
A map

\[
f:X\to Y
\]

is a \emph{homeomorphism} if:

\begin{enumerate}
    \item \(f\) is bijective;
    \item \(f\) is continuous;
    \item \(f^{-1}\) is continuous.
\end{enumerate}
\end{definitionbox}

If a homeomorphism exists, \(X\) and \(Y\) are called
\emph{homeomorphic}.

We write

\[
\boxed{
X\cong Y.
}
\]

%%%%%%%%%%%%%%%%%%%%%%%%%%%%%%%%%%%%%%%%%%%%%%%%%%%%%%%%%%%%
\subsection{Topological Invariants}
\label{subsec:topological-invariants}
%%%%%%%%%%%%%%%%%%%%%%%%%%%%%%%%%%%%%%%%%%%%%%%%%%%%%%%%%%%%

A property preserved by homeomorphisms is called a
\emph{topological invariant}.

Examples include:

\[
\boxed{
\text{compactness},
\quad
\text{connectedness},
\quad
\text{number of connected components}.
}
\]

Later sections develop more sophisticated invariants such as

\[
\boxed{
\text{fundamental groups}
}
\]

and

\[
\boxed{
\text{homology groups}.
}
\]

%%%%%%%%%%%%%%%%%%%%%%%%%%%%%%%%%%%%%%%%%%%%%%%%%%%%%%%%%%%%
\subsection{Sequences in Topological Spaces}
\label{subsec:topological-sequences}
%%%%%%%%%%%%%%%%%%%%%%%%%%%%%%%%%%%%%%%%%%%%%%%%%%%%%%%%%%%%

Let

\[
(x_n)
\]

be a sequence in a topological space \(X\).

\begin{definitionbox}{Sequence Convergence}
The sequence \((x_n)\) \emph{converges} to \(x\in X\) if for every
neighborhood \(U\) of \(x\), there exists \(N\in\mathbb{N}\) such that

\[
\boxed{
n\geq N
\Longrightarrow
x_n\in U.
}
\]
\end{definitionbox}

We write

\[
\boxed{
x_n\to x.
}
\]

%%%%%%%%%%%%%%%%%%%%%%%%%%%%%%%%%%%%%%%%%%%%%%%%%%%%%%%%%%%%
\subsection{Metric and Topological Convergence}
\label{subsec:metric-topological-convergence}
%%%%%%%%%%%%%%%%%%%%%%%%%%%%%%%%%%%%%%%%%%%%%%%%%%%%%%%%%%%%

In a metric space,

\[
x_n\to x
\]

means

\[
\boxed{
\forall\varepsilon>0,\
\exists N,\
n\geq N
\Longrightarrow
d(x_n,x)<\varepsilon.
}
\]

The lowercase Greek letter

\[
\varepsilon
\]

is read ``epsilon.''

This is equivalent to saying that the sequence eventually lies inside every
neighborhood of \(x\).

Thus the topological definition generalizes metric convergence.

%%%%%%%%%%%%%%%%%%%%%%%%%%%%%%%%%%%%%%%%%%%%%%%%%%%%%%%%%%%%
\subsection{Sequences Do Not Always Detect Everything}
\label{subsec:sequences-not-enough}
%%%%%%%%%%%%%%%%%%%%%%%%%%%%%%%%%%%%%%%%%%%%%%%%%%%%%%%%%%%%

In metric spaces, sequences are powerful enough to characterize many
topological properties.

In arbitrary topological spaces, this is no longer always true.

For example, closure cannot always be completely characterized merely by
limits of sequences.

More general objects called \emph{nets} and \emph{filters} can recover the
full topology.

These will be treated later when needed.

%%%%%%%%%%%%%%%%%%%%%%%%%%%%%%%%%%%%%%%%%%%%%%%%%%%%%%%%%%%%
\subsection{Hausdorff Spaces}
\label{subsec:hausdorff-spaces}
%%%%%%%%%%%%%%%%%%%%%%%%%%%%%%%%%%%%%%%%%%%%%%%%%%%%%%%%%%%%

\begin{definitionbox}{Hausdorff Space}
A topological space \(X\) is \emph{Hausdorff} if for every pair of distinct
points \(x,y\in X\), there exist disjoint open sets \(U,V\) such that

\[
\boxed{
x\in U,
\qquad
y\in V,
\qquad
U\cap V=\varnothing.
}
\]
\end{definitionbox}

Hausdorff spaces are also called \(T_2\) spaces.

%%%%%%%%%%%%%%%%%%%%%%%%%%%%%%%%%%%%%%%%%%%%%%%%%%%%%%%%%%%%
\subsection{Why Hausdorff Matters}
\label{subsec:why-hausdorff-matters}
%%%%%%%%%%%%%%%%%%%%%%%%%%%%%%%%%%%%%%%%%%%%%%%%%%%%%%%%%%%%

In a Hausdorff space, limits of sequences are unique.

\begin{theorem}[Uniqueness of Limits]
If \(X\) is Hausdorff and

\[
x_n\to x
\]

and

\[
x_n\to y,
\]

then

\[
\boxed{x=y.}
\]
\end{theorem}

This is one reason the Hausdorff condition appears throughout analysis and
geometry.

%%%%%%%%%%%%%%%%%%%%%%%%%%%%%%%%%%%%%%%%%%%%%%%%%%%%%%%%%%%%
\subsection{Metric Spaces Are Hausdorff}
\label{subsec:metric-spaces-hausdorff}
%%%%%%%%%%%%%%%%%%%%%%%%%%%%%%%%%%%%%%%%%%%%%%%%%%%%%%%%%%%%

Suppose

\[
x\neq y
\]

in a metric space.

Then

\[
d(x,y)>0.
\]

Choose

\[
r=\frac{d(x,y)}{3}.
\]

The open balls

\[
B_r(x)
\]

and

\[
B_r(y)
\]

are disjoint.

Therefore

\[
\boxed{
\text{every metric space is Hausdorff}.
}
\]

%%%%%%%%%%%%%%%%%%%%%%%%%%%%%%%%%%%%%%%%%%%%%%%%%%%%%%%%%%%%
\subsection{Separation Axioms}
\label{subsec:separation-axioms}
%%%%%%%%%%%%%%%%%%%%%%%%%%%%%%%%%%%%%%%%%%%%%%%%%%%%%%%%%%%%

Topology contains a hierarchy of conditions describing how effectively
points and closed sets can be separated by open sets.

Common conditions include

\[
\boxed{
T_0,
\quad
T_1,
\quad
T_2,
\quad
T_3,
\quad
T_4.
}
\]

The symbol \(T\) comes from the German word \emph{Trennungsaxiom}, meaning
``separation axiom.''

Conventions for \(T_3\) and \(T_4\) can vary slightly between texts.

%%%%%%%%%%%%%%%%%%%%%%%%%%%%%%%%%%%%%%%%%%%%%%%%%%%%%%%%%%%%
\subsection{\texorpdfstring{\(T_0\)}{T0} Spaces}
\label{subsec:t0-spaces}
%%%%%%%%%%%%%%%%%%%%%%%%%%%%%%%%%%%%%%%%%%%%%%%%%%%%%%%%%%%%

A space is \(T_0\) if any two distinct points are topologically
distinguishable.

That is, for distinct \(x,y\), there is an open set containing one but not
the other.

This is the weakest of the standard separation conditions considered here.

%%%%%%%%%%%%%%%%%%%%%%%%%%%%%%%%%%%%%%%%%%%%%%%%%%%%%%%%%%%%
\subsection{\texorpdfstring{\(T_1\)}{T1} Spaces}
\label{subsec:t1-spaces}
%%%%%%%%%%%%%%%%%%%%%%%%%%%%%%%%%%%%%%%%%%%%%%%%%%%%%%%%%%%%

A space is \(T_1\) if for any two distinct points \(x,y\), each has an open
neighborhood not containing the other.

An equivalent condition is:

\[
\boxed{
\{x\}\text{ is closed for every }x\in X.
}
\]

Every Hausdorff space is \(T_1\).

%%%%%%%%%%%%%%%%%%%%%%%%%%%%%%%%%%%%%%%%%%%%%%%%%%%%%%%%%%%%
\subsection{Compactness}
\label{subsec:topological-compactness}
%%%%%%%%%%%%%%%%%%%%%%%%%%%%%%%%%%%%%%%%%%%%%%%%%%%%%%%%%%%%

Compactness generalizes the finite character of closed and bounded intervals.

\begin{definitionbox}{Open Cover}
An \emph{open cover} of \(X\) is a collection of open sets

\[
\{U_\alpha\}_{\alpha\in A}
\]

such that

\[
\boxed{
X
\subseteq
\bigcup_{\alpha\in A}U_\alpha.
}
\]
\end{definitionbox}

%%%%%%%%%%%%%%%%%%%%%%%%%%%%%%%%%%%%%%%%%%%%%%%%%%%%%%%%%%%%
\subsection{Compact Spaces}
\label{subsec:compact-spaces}
%%%%%%%%%%%%%%%%%%%%%%%%%%%%%%%%%%%%%%%%%%%%%%%%%%%%%%%%%%%%

\begin{definitionbox}{Compact Space}
A topological space \(X\) is \emph{compact} if every open cover of \(X\)
contains a finite subcover.
\end{definitionbox}

That is, whenever

\[
X\subseteq
\bigcup_{\alpha\in A}U_\alpha,
\]

there exist

\[
\alpha_1,\ldots,\alpha_n
\]

such that

\[
\boxed{
X
\subseteq
U_{\alpha_1}\cup\cdots\cup U_{\alpha_n}.
}
\]

%%%%%%%%%%%%%%%%%%%%%%%%%%%%%%%%%%%%%%%%%%%%%%%%%%%%%%%%%%%%
\subsection{Compact Subsets}
\label{subsec:compact-subsets}
%%%%%%%%%%%%%%%%%%%%%%%%%%%%%%%%%%%%%%%%%%%%%%%%%%%%%%%%%%%%

A subset \(K\subseteq X\) is compact when it is compact with respect to its
subspace topology.

Equivalently, every collection of open subsets of \(X\) covering \(K\)
contains a finite subcollection that still covers \(K\).

%%%%%%%%%%%%%%%%%%%%%%%%%%%%%%%%%%%%%%%%%%%%%%%%%%%%%%%%%%%%
\subsection{Heine--Borel Theorem}
\label{subsec:heine-borel-topology}
%%%%%%%%%%%%%%%%%%%%%%%%%%%%%%%%%%%%%%%%%%%%%%%%%%%%%%%%%%%%

In Euclidean space, compactness has a familiar characterization.

\begin{theorem}[Heine--Borel Theorem]
A subset

\[
K\subseteq\R^n
\]

is compact if and only if it is closed and bounded.
\end{theorem}

\begin{warningbox}
``Compact'' does not mean ``closed and bounded'' in every topological space.

The closed-and-bounded characterization is a special property of
finite-dimensional Euclidean spaces and related settings.
\end{warningbox}

%%%%%%%%%%%%%%%%%%%%%%%%%%%%%%%%%%%%%%%%%%%%%%%%%%%%%%%%%%%%
\subsection{Continuous Images of Compact Spaces}
\label{subsec:continuous-image-compact}
%%%%%%%%%%%%%%%%%%%%%%%%%%%%%%%%%%%%%%%%%%%%%%%%%%%%%%%%%%%%

\begin{theorem}
Let

\[
f:X\to Y
\]

be continuous.

If \(X\) is compact, then

\[
\boxed{
f(X)
\text{ is compact}.
}
\]
\end{theorem}

This theorem is purely topological.

It does not require a metric.

%%%%%%%%%%%%%%%%%%%%%%%%%%%%%%%%%%%%%%%%%%%%%%%%%%%%%%%%%%%%
\subsection{Compact Sets in Hausdorff Spaces}
\label{subsec:compact-hausdorff}
%%%%%%%%%%%%%%%%%%%%%%%%%%%%%%%%%%%%%%%%%%%%%%%%%%%%%%%%%%%%

\begin{theorem}
Every compact subset of a Hausdorff space is closed.
\end{theorem}

This theorem connects compactness with the separation properties of the
ambient space.

%%%%%%%%%%%%%%%%%%%%%%%%%%%%%%%%%%%%%%%%%%%%%%%%%%%%%%%%%%%%
\subsection{Extreme Value Theorem as a Topological Consequence}
\label{subsec:extreme-value-topological}
%%%%%%%%%%%%%%%%%%%%%%%%%%%%%%%%%%%%%%%%%%%%%%%%%%%%%%%%%%%%

Suppose

\[
f:K\to\R
\]

is continuous and \(K\) is compact.

Then \(f(K)\) is compact.

Since compact subsets of \(\R\) are closed and bounded, \(f(K)\) contains
its maximum and minimum.

Thus the Extreme Value Theorem can be understood through compactness.

%%%%%%%%%%%%%%%%%%%%%%%%%%%%%%%%%%%%%%%%%%%%%%%%%%%%%%%%%%%%
\subsection{Connectedness}
\label{subsec:topological-connectedness}
%%%%%%%%%%%%%%%%%%%%%%%%%%%%%%%%%%%%%%%%%%%%%%%%%%%%%%%%%%%%

\begin{definitionbox}{Connected Space}
A topological space \(X\) is \emph{connected} if it cannot be written as

\[
\boxed{
X=U\cup V
}
\]

where \(U\) and \(V\) are nonempty, disjoint open subsets of \(X\).
\end{definitionbox}

Such a decomposition, when it exists, is called a \emph{separation} of
\(X\).

%%%%%%%%%%%%%%%%%%%%%%%%%%%%%%%%%%%%%%%%%%%%%%%%%%%%%%%%%%%%
\subsection{Connectedness and Clopen Sets}
\label{subsec:connected-clopen}
%%%%%%%%%%%%%%%%%%%%%%%%%%%%%%%%%%%%%%%%%%%%%%%%%%%%%%%%%%%%

A space \(X\) is connected if and only if its only clopen subsets are

\[
\boxed{
\varnothing
\quad\text{and}\quad
X.
}
\]

Thus nontrivial clopen sets reveal disconnectedness.

%%%%%%%%%%%%%%%%%%%%%%%%%%%%%%%%%%%%%%%%%%%%%%%%%%%%%%%%%%%%
\subsection{Intervals Are Connected}
\label{subsec:intervals-connected}
%%%%%%%%%%%%%%%%%%%%%%%%%%%%%%%%%%%%%%%%%%%%%%%%%%%%%%%%%%%%

Every interval in \(\R\) is connected.

This includes

\[
(a,b),
\qquad
[a,b],
\qquad
(a,b],
\qquad
[a,b),
\]

as well as unbounded intervals.

In fact, the connected subsets of \(\R\) are precisely the intervals,
including singletons and the empty set under the usual conventions.

%%%%%%%%%%%%%%%%%%%%%%%%%%%%%%%%%%%%%%%%%%%%%%%%%%%%%%%%%%%%
\subsection{Continuous Images Preserve Connectedness}
\label{subsec:continuous-image-connected}
%%%%%%%%%%%%%%%%%%%%%%%%%%%%%%%%%%%%%%%%%%%%%%%%%%%%%%%%%%%%

\begin{theorem}
If \(X\) is connected and

\[
f:X\to Y
\]

is continuous, then

\[
\boxed{
f(X)
\text{ is connected}.
}
\]
\end{theorem}

This theorem provides a topological explanation for the Intermediate Value
Theorem.

%%%%%%%%%%%%%%%%%%%%%%%%%%%%%%%%%%%%%%%%%%%%%%%%%%%%%%%%%%%%
\subsection{Intermediate Value Theorem and Connectedness}
\label{subsec:ivt-connectedness}
%%%%%%%%%%%%%%%%%%%%%%%%%%%%%%%%%%%%%%%%%%%%%%%%%%%%%%%%%%%%

Let

\[
f:[a,b]\to\R
\]

be continuous.

Since

\[
[a,b]
\]

is connected, its image

\[
f([a,b])
\]

must also be connected.

Connected subsets of \(\R\) are intervals.

Therefore \(f([a,b])\) contains every value between any two of its values.

This is precisely the geometric content behind the Intermediate Value
Theorem.

%%%%%%%%%%%%%%%%%%%%%%%%%%%%%%%%%%%%%%%%%%%%%%%%%%%%%%%%%%%%
\subsection{Path Connectedness}
\label{subsec:path-connectedness}
%%%%%%%%%%%%%%%%%%%%%%%%%%%%%%%%%%%%%%%%%%%%%%%%%%%%%%%%%%%%

\begin{definitionbox}{Path}
A \emph{path} from \(x\) to \(y\) in \(X\) is a continuous function

\[
\boxed{
\gamma:[0,1]\to X
}
\]

such that

\[
\gamma(0)=x
\]

and

\[
\gamma(1)=y.
\]
\end{definitionbox}

The lowercase Greek letter

\[
\gamma
\]

is read ``gamma.''

%%%%%%%%%%%%%%%%%%%%%%%%%%%%%%%%%%%%%%%%%%%%%%%%%%%%%%%%%%%%
\subsection{Path-Connected Spaces}
\label{subsec:path-connected-spaces}
%%%%%%%%%%%%%%%%%%%%%%%%%%%%%%%%%%%%%%%%%%%%%%%%%%%%%%%%%%%%

\begin{definitionbox}{Path Connected}
A space \(X\) is \emph{path connected} if every pair of points
\(x,y\in X\) can be joined by a path in \(X\).
\end{definitionbox}

Every path-connected space is connected:

\[
\boxed{
\text{path connected}
\Longrightarrow
\text{connected}.
}
\]

The converse is not true in general.

%%%%%%%%%%%%%%%%%%%%%%%%%%%%%%%%%%%%%%%%%%%%%%%%%%%%%%%%%%%%
\subsection{Connected Components}
\label{subsec:connected-components}
%%%%%%%%%%%%%%%%%%%%%%%%%%%%%%%%%%%%%%%%%%%%%%%%%%%%%%%%%%%%

A \emph{connected component} of \(X\) is a maximal connected subset of \(X\).

The connected components partition the space.

Similarly, path components are maximal path-connected subsets.

These notions need not coincide in arbitrary spaces.

%%%%%%%%%%%%%%%%%%%%%%%%%%%%%%%%%%%%%%%%%%%%%%%%%%%%%%%%%%%%
\subsection{Homeomorphism Preserves Connectedness}
\label{subsec:homeomorphism-connectedness}
%%%%%%%%%%%%%%%%%%%%%%%%%%%%%%%%%%%%%%%%%%%%%%%%%%%%%%%%%%%%

If

\[
X\cong Y,
\]

then

\[
\boxed{
X\text{ connected}
\iff
Y\text{ connected}.
}
\]

Likewise, compactness and path-connectedness are preserved by
homeomorphisms.

These properties can therefore be used to prove that spaces are not
homeomorphic.

%%%%%%%%%%%%%%%%%%%%%%%%%%%%%%%%%%%%%%%%%%%%%%%%%%%%%%%%%%%%
\subsection{Worked Example: Removing a Point}
\label{subsec:worked-removing-point}
%%%%%%%%%%%%%%%%%%%%%%%%%%%%%%%%%%%%%%%%%%%%%%%%%%%%%%%%%%%%

\begin{workedexamplebox}{Distinguishing the Line and Plane}

Consider

\[
\R
\]

and

\[
\R^2.
\]

Remove the origin.

Then

\[
\R\setminus\{0\}
=
(-\infty,0)\cup(0,\infty)
\]

is disconnected.

However,

\[
\R^2\setminus\{(0,0)\}
\]

remains path connected: one can move around the missing point.

Therefore \(\R\) and \(\R^2\) cannot be homeomorphic.

\begin{answerbox}
\[
\boxed{
\R\not\cong\R^2.
}
\]
\end{answerbox}
\end{workedexamplebox}

%%%%%%%%%%%%%%%%%%%%%%%%%%%%%%%%%%%%%%%%%%%%%%%%%%%%%%%%%%%%
\subsection{Topological Properties Versus Geometric Properties}
\label{subsec:topological-vs-geometric-properties}
%%%%%%%%%%%%%%%%%%%%%%%%%%%%%%%%%%%%%%%%%%%%%%%%%%%%%%%%%%%%

Examples of topological properties include

\[
\boxed{
\text{compactness},
\quad
\text{connectedness},
\quad
\text{Hausdorffness}.
}
\]

Examples of properties that are generally not topological include

\[
\boxed{
\text{length},
\quad
\text{angle},
\quad
\text{area},
\quad
\text{curvature}.
}
\]

A homeomorphism may distort all of these geometric quantities.

%%%%%%%%%%%%%%%%%%%%%%%%%%%%%%%%%%%%%%%%%%%%%%%%%%%%%%%%%%%%
\subsection{Compactness and Connectedness of Familiar Spaces}
\label{subsec:familiar-space-properties}
%%%%%%%%%%%%%%%%%%%%%%%%%%%%%%%%%%%%%%%%%%%%%%%%%%%%%%%%%%%%

In the usual topology:

\[
[0,1]
\]

is compact and connected.

The interval

\[
(0,1)
\]

is connected but not compact.

The set

\[
\Z
\]

is closed but disconnected.

The space

\[
\R
\]

is connected but not compact.

The circle

\[
S^1
\]

is compact and connected.

%%%%%%%%%%%%%%%%%%%%%%%%%%%%%%%%%%%%%%%%%%%%%%%%%%%%%%%%%%%%
\subsection{Continuous Bijections Need Not Be Homeomorphisms}
\label{subsec:continuous-bijection-not-homeomorphism}
%%%%%%%%%%%%%%%%%%%%%%%%%%%%%%%%%%%%%%%%%%%%%%%%%%%%%%%%%%%%

A continuous bijection

\[
f:X\to Y
\]

does not automatically have a continuous inverse.

Therefore continuity and bijectivity alone do not always establish a
homeomorphism.

However, an important theorem gives a useful sufficient condition.

\begin{theorem}
If \(X\) is compact, \(Y\) is Hausdorff, and

\[
f:X\to Y
\]

is a continuous bijection, then \(f\) is a homeomorphism.
\end{theorem}

%%%%%%%%%%%%%%%%%%%%%%%%%%%%%%%%%%%%%%%%%%%%%%%%%%%%%%%%%%%%
\subsection{Compactness as a Finiteness Principle}
\label{subsec:compactness-finiteness-principle}
%%%%%%%%%%%%%%%%%%%%%%%%%%%%%%%%%%%%%%%%%%%%%%%%%%%%%%%%%%%%

Compactness often allows an infinite collection of local information to be
reduced to finitely many pieces.

The defining pattern is

\[
\boxed{
\text{arbitrary open cover}
\Longrightarrow
\text{finite subcover}.
}
\]

This principle appears repeatedly throughout analysis, geometry, and
differential equations.

%%%%%%%%%%%%%%%%%%%%%%%%%%%%%%%%%%%%%%%%%%%%%%%%%%%%%%%%%%%%
\subsection{Local and Global Properties}
\label{subsec:local-global-topology}
%%%%%%%%%%%%%%%%%%%%%%%%%%%%%%%%%%%%%%%%%%%%%%%%%%%%%%%%%%%%

Some topological properties are local.

For example, a manifold is locally modeled on Euclidean space.

Other properties are global.

Compactness and connectedness describe the space as a whole.

This local-global distinction becomes central in later topology and
differential geometry.

%%%%%%%%%%%%%%%%%%%%%%%%%%%%%%%%%%%%%%%%%%%%%%%%%%%%%%%%%%%%
\subsection{Common Errors}
\label{subsec:point-set-topology-common-errors}
%%%%%%%%%%%%%%%%%%%%%%%%%%%%%%%%%%%%%%%%%%%%%%%%%%%%%%%%%%%%

\subsubsection{Assuming Open Means ``Does Not Include Endpoints''}

Openness is defined relative to a topology.

The phrase ``does not include endpoints'' applies only to certain familiar
intervals in \(\R\).

A general topological space may have no meaningful notion of endpoint.

%%%%%%%%%%%%%%%%%%%%%%%%%%%%%%%%%%%%%%%%%%%%%%%%%%%%%%%%%%%%
\subsubsection{Assuming Closed Means ``Contains Endpoints''}

A set is closed when its complement is open.

Endpoint language is not the general definition.

%%%%%%%%%%%%%%%%%%%%%%%%%%%%%%%%%%%%%%%%%%%%%%%%%%%%%%%%%%%%
\subsubsection{Assuming Open and Closed Are Opposites}

A set may be both open and closed.

For example,

\[
\varnothing
\]

and

\[
X
\]

are always both.

%%%%%%%%%%%%%%%%%%%%%%%%%%%%%%%%%%%%%%%%%%%%%%%%%%%%%%%%%%%%
\subsubsection{Taking Infinite Intersections of Open Sets}

Finite intersections of open sets are guaranteed to be open.

Infinite intersections need not be open.

%%%%%%%%%%%%%%%%%%%%%%%%%%%%%%%%%%%%%%%%%%%%%%%%%%%%%%%%%%%%
\subsubsection{Taking Infinite Unions of Closed Sets}

Finite unions of closed sets are closed.

Infinite unions need not be closed.

For example,

\[
\bigcup_{n=1}^{\infty}
\left[
\frac1n,1
\right]
=
(0,1]
\]

is not closed in \(\R\).

%%%%%%%%%%%%%%%%%%%%%%%%%%%%%%%%%%%%%%%%%%%%%%%%%%%%%%%%%%%%
\subsubsection{Confusing Interior with the Original Set}

In general,

\[
\operatorname{int}(A)
\subseteq A,
\]

and equality occurs precisely when \(A\) is open.

%%%%%%%%%%%%%%%%%%%%%%%%%%%%%%%%%%%%%%%%%%%%%%%%%%%%%%%%%%%%
\subsubsection{Confusing Closure with Boundary}

The closure contains the original set:

\[
A\subseteq\overline{A}.
\]

The boundary is only

\[
\partial A
=
\overline{A}\setminus\operatorname{int}(A).
\]

%%%%%%%%%%%%%%%%%%%%%%%%%%%%%%%%%%%%%%%%%%%%%%%%%%%%%%%%%%%%
\subsubsection{Forgetting the Ambient Space}

Interior, closure, and boundary depend on the ambient topology.

For example,

\[
[0,1]
\]

has different relative interior when considered as a subspace of itself than
when considered as a subset of \(\R\).

%%%%%%%%%%%%%%%%%%%%%%%%%%%%%%%%%%%%%%%%%%%%%%%%%%%%%%%%%%%%
\subsubsection{Confusing Limit Points with Points of the Set}

A limit point need not belong to the set.

For example,

\[
0
\]

is a limit point of

\[
(0,1)
\]

but

\[
0\notin(0,1).
\]

%%%%%%%%%%%%%%%%%%%%%%%%%%%%%%%%%%%%%%%%%%%%%%%%%%%%%%%%%%%%
\subsubsection{Assuming Dense Means Equal to the Whole Space}

The rational numbers satisfy

\[
\Q\neq\R
\]

but

\[
\overline{\Q}=\R.
\]

Thus density concerns closure, not equality.

%%%%%%%%%%%%%%%%%%%%%%%%%%%%%%%%%%%%%%%%%%%%%%%%%%%%%%%%%%%%
\subsubsection{Confusing Basis Elements with All Open Sets}

A basis need not contain every open set.

Open sets are arbitrary unions of basis elements.

%%%%%%%%%%%%%%%%%%%%%%%%%%%%%%%%%%%%%%%%%%%%%%%%%%%%%%%%%%%%
\subsubsection{Confusing Open in a Subspace with Open in the Ambient Space}

Relative openness depends on the subspace topology.

A set may be open in \(Y\) while not being open in \(X\).

%%%%%%%%%%%%%%%%%%%%%%%%%%%%%%%%%%%%%%%%%%%%%%%%%%%%%%%%%%%%
\subsubsection{Using Images Instead of Preimages for Continuity}

Topological continuity requires

\[
\boxed{
f^{-1}(U)\text{ open whenever }U\text{ is open}.
}
\]

A continuous function need not send open sets to open sets.

%%%%%%%%%%%%%%%%%%%%%%%%%%%%%%%%%%%%%%%%%%%%%%%%%%%%%%%%%%%%
\subsubsection{Assuming a Continuous Bijection Is a Homeomorphism}

The inverse must also be continuous.

Additional hypotheses, such as compact domain and Hausdorff codomain, can
guarantee this.

%%%%%%%%%%%%%%%%%%%%%%%%%%%%%%%%%%%%%%%%%%%%%%%%%%%%%%%%%%%%
\subsubsection{Assuming Every Space Is Hausdorff}

Hausdorffness is an additional property, not part of the definition of a
topological space.

%%%%%%%%%%%%%%%%%%%%%%%%%%%%%%%%%%%%%%%%%%%%%%%%%%%%%%%%%%%%
\subsubsection{Assuming Compact Means Closed and Bounded Everywhere}

The Heine--Borel characterization is specific to Euclidean spaces and
certain related settings.

Compactness itself is defined using open covers.

%%%%%%%%%%%%%%%%%%%%%%%%%%%%%%%%%%%%%%%%%%%%%%%%%%%%%%%%%%%%
\subsubsection{Confusing Connected and Path Connected}

Path connectedness implies connectedness.

The converse fails in general.

%%%%%%%%%%%%%%%%%%%%%%%%%%%%%%%%%%%%%%%%%%%%%%%%%%%%%%%%%%%%
\subsubsection{Assuming Sequences Capture Every Topological Property}

Sequences work especially well in metric and first-countable spaces.

General topological spaces may require nets or filters.

%%%%%%%%%%%%%%%%%%%%%%%%%%%%%%%%%%%%%%%%%%%%%%%%%%%%%%%%%%%%
\subsection{Applications and Connections}
\label{subsec:point-set-topology-applications}
%%%%%%%%%%%%%%%%%%%%%%%%%%%%%%%%%%%%%%%%%%%%%%%%%%%%%%%%%%%%

\begin{applicationbox}
\textbf{Analysis.}
Open sets, neighborhoods, compactness, connectedness, and continuity provide
the abstract framework underlying real and complex analysis.

\medskip

\textbf{Metric Spaces.}
Every metric induces a topology, allowing metric notions of convergence and
continuity to be expressed topologically.

\medskip

\textbf{Geometry.}
Topology distinguishes geometric structure that survives continuous
deformation from quantities such as length, angle, and curvature.

\medskip

\textbf{Differential Geometry.}
Smooth manifolds begin as topological spaces locally modeled on Euclidean
space before differentiable structure is added.

\medskip

\textbf{Algebraic Topology.}
Homeomorphisms, connectedness, quotient spaces, and paths form the basic
language needed to define homotopy, fundamental groups, and homology.

\medskip

\textbf{Functional Analysis.}
Infinite-dimensional spaces are studied using metric, norm, weak, and other
topologies.

\medskip

\textbf{Differential Equations.}
Compactness and continuity are fundamental tools in existence arguments and
qualitative analysis.

\medskip

\textbf{Optimization.}
Compactness often guarantees that continuous objective functions attain
maximum and minimum values.

\medskip

\textbf{Data Analysis.}
Topological data analysis studies qualitative shape and connectivity in
data using topological invariants.
\end{applicationbox}

%%%%%%%%%%%%%%%%%%%%%%%%%%%%%%%%%%%%%%%%%%%%%%%%%%%%%%%%%%%%
\subsection{Exercises}
\label{subsec:point-set-topology-exercises}
%%%%%%%%%%%%%%%%%%%%%%%%%%%%%%%%%%%%%%%%%%%%%%%%%%%%%%%%%%%%

\begin{exercise}[1]
State the three axioms defining a topology.
\end{exercise}

\begin{exercise}[1]
What is the indiscrete topology on a set \(X\)?
\end{exercise}

\begin{exercise}[1]
What is the discrete topology on \(X\)?
\end{exercise}

\begin{exercise}[1]
Which topology on \(X\) is the finest possible topology?
\end{exercise}

\begin{exercise}[1]
Which topology on \(X\) is the coarsest possible topology?
\end{exercise}

\begin{exercise}[1]
Explain why arbitrary unions of open sets are open.
\end{exercise}

\begin{exercise}[1]
Give an example showing that an infinite intersection of open subsets of
\(\R\) need not be open.
\end{exercise}

\begin{exercise}[1]
Define a closed subset of a topological space.
\end{exercise}

\begin{exercise}[1]
Can a subset be both open and closed?
\end{exercise}

\begin{exercise}[1]
Define a neighborhood of a point.
\end{exercise}

\begin{exercise}[1]
Define the interior of \(A\subseteq X\).
\end{exercise}

\begin{exercise}[1]
Define the closure of \(A\subseteq X\).
\end{exercise}

\begin{exercise}[1]
Define the boundary of \(A\subseteq X\).
\end{exercise}

\begin{exercise}[1]
For

\[
A=(2,5]
\subseteq\R,
\]

find

\[
\operatorname{int}(A),
\qquad
\overline{A},
\qquad
\partial A.
\]
\end{exercise}

\begin{exercise}[1]
Define a limit point.
\end{exercise}

\begin{exercise}[1]
Define an isolated point.
\end{exercise}

\begin{exercise}[1]
What does it mean for \(A\) to be dense in \(X\)?
\end{exercise}

\begin{exercise}[1]
Give a proper dense subset of \(\R\).
\end{exercise}

\begin{exercise}[1]
Define a basis for a topology.
\end{exercise}

\begin{exercise}[1]
What is the standard interval basis for \(\R\)?
\end{exercise}

\begin{exercise}[1]
Define the subspace topology.
\end{exercise}

\begin{exercise}[1]
Define the product topology on \(X\times Y\).
\end{exercise}

\begin{exercise}[1]
Define topological continuity.
\end{exercise}

\begin{exercise}[1]
Define a homeomorphism.
\end{exercise}

\begin{exercise}[1]
Define compactness.
\end{exercise}

\begin{exercise}[1]
Define connectedness.
\end{exercise}

\begin{exercise}[1]
Define path connectedness.
\end{exercise}

\begin{exercise}[1]
Define a Hausdorff space.
\end{exercise}

\begin{exercise}[2]
Let

\[
X=\{1,2,3\}.
\]

Determine whether

\[
\mathcal{T}
=
\{
\varnothing,
X,
\{1\},
\{1,2\}
\}
\]

is a topology.
\end{exercise}

\begin{exercise}[2]
Determine whether

\[
\mathcal{S}
=
\{
\varnothing,
X,
\{1\},
\{2\}
\}
\]

is a topology on the same set.
\end{exercise}

\begin{exercise}[2]
List every open set in the discrete topology on

\[
X=\{a,b\}.
\]
\end{exercise}

\begin{exercise}[2]
Determine the closed sets in the indiscrete topology on a nonempty set
\(X\).
\end{exercise}

\begin{exercise}[2]
Show that arbitrary intersections of closed sets are closed.
\end{exercise}

\begin{exercise}[2]
Show that finite unions of closed sets are closed.
\end{exercise}

\begin{exercise}[2]
For

\[
A=\Q\subseteq\R,
\]

determine its interior and closure.
\end{exercise}

\begin{exercise}[2]
For

\[
A=\Z\subseteq\R,
\]

determine its interior, closure, and set of limit points.
\end{exercise}

\begin{exercise}[2]
Find the boundary of

\[
[0,1]\cup[2,3]
\]

in \(\R\).
\end{exercise}

\begin{exercise}[2]
Show that

\[
\operatorname{int}(A)
\]

is the largest open subset contained in \(A\).
\end{exercise}

\begin{exercise}[2]
Show that

\[
\overline{A}
\]

is the smallest closed subset containing \(A\).
\end{exercise}

\begin{exercise}[2]
Show that

\[
\partial A
=
\overline{A}\cap\overline{X\setminus A}.
\]
\end{exercise}

\begin{exercise}[2]
Let

\[
Y=[0,1]\subseteq\R.
\]

Show that

\[
[0,1/2)
\]

is open in \(Y\).
\end{exercise}

\begin{exercise}[2]
Give a subset of \(Y=[0,1]\) that is closed in \(Y\) but not closed in
\(\R\).
\end{exercise}

\begin{exercise}[2]
Explain why open rectangles form a basis for the usual topology on
\(\R^2\).
\end{exercise}

\begin{exercise}[2]
Show that the composition of two continuous maps is continuous.
\end{exercise}

\begin{exercise}[2]
Show that every function from a discrete space into any topological space is
continuous.
\end{exercise}

\begin{exercise}[2]
Determine which functions from an indiscrete space to an arbitrary
topological space are necessarily continuous.
\end{exercise}

\begin{exercise}[2]
Explain why every metric space is Hausdorff.
\end{exercise}

\begin{exercise}[2]
Determine whether

\[
(0,1)
\]

is compact in \(\R\).
\end{exercise}

\begin{exercise}[2]
Determine whether

\[
[0,1]
\]

is compact in \(\R\).
\end{exercise}

\begin{exercise}[2]
Determine whether

\[
\R
\]

is connected.
\end{exercise}

\begin{exercise}[2]
Determine whether

\[
\R\setminus\{0\}
\]

is connected.
\end{exercise}

\begin{exercise}[2]
Show that every convex subset of \(\R^n\) is path connected.
\end{exercise}

\begin{exercise}[3]
Prove that

\[
\overline{A\cup B}
=
\overline{A}\cup\overline{B}.
\]
\end{exercise}

\begin{exercise}[3]
Investigate whether

\[
\overline{A\cap B}
=
\overline{A}\cap\overline{B}
\]

always holds.
\end{exercise}

\begin{exercise}[3]
Prove that

\[
\operatorname{int}(A\cap B)
=
\operatorname{int}(A)\cap\operatorname{int}(B).
\]
\end{exercise}

\begin{exercise}[3]
Investigate whether

\[
\operatorname{int}(A\cup B)
=
\operatorname{int}(A)\cup\operatorname{int}(B)
\]

always holds.
\end{exercise}

\begin{exercise}[3]
Prove that

\[
x\in\overline{A}
\]

if and only if every open neighborhood of \(x\) intersects \(A\).
\end{exercise}

\begin{exercise}[3]
Prove that

\[
A\text{ is closed}
\iff
A
\text{ contains all of its limit points}.
\]
\end{exercise}

\begin{exercise}[3]
Show that intervals with rational endpoints form a countable basis for the
usual topology on \(\R\).
\end{exercise}

\begin{exercise}[3]
Prove that every second-countable space is separable under the usual mild
choice assumptions.
\end{exercise}

\begin{exercise}[3]
Show that the quotient of

\[
[0,1]
\]

obtained by identifying \(0\) and \(1\) is homeomorphic to \(S^1\).
\end{exercise}

\begin{exercise}[3]
Prove that a continuous image of a compact space is compact.
\end{exercise}

\begin{exercise}[3]
Prove that a compact subset of a Hausdorff space is closed.
\end{exercise}

\begin{exercise}[3]
Prove that a continuous image of a connected space is connected.
\end{exercise}

\begin{exercise}[3]
Prove that every path-connected space is connected.
\end{exercise}

\begin{exercise}[3]
Give an example of a connected space that is not path connected.
\end{exercise}

\begin{exercise}[3]
Prove that limits of sequences are unique in Hausdorff spaces.
\end{exercise}

\begin{exercise}[3]
Prove that a continuous bijection from a compact space to a Hausdorff space
is a homeomorphism.
\end{exercise}

\begin{exercise}[3]
Use connectedness to prove the Intermediate Value Theorem.
\end{exercise}

\begin{exercise}[3]
Use compactness to explain the topological structure behind the Extreme
Value Theorem.
\end{exercise}

\begin{exercise}[4]
Construct a topology on a three-element set that is neither discrete nor
indiscrete.
\end{exercise}

\begin{exercise}[4]
Find all topologies on a two-element set.
\end{exercise}

\begin{exercise}[4]
Investigate the cofinite topology on an infinite set and determine whether
it is \(T_1\) and whether it is Hausdorff.
\end{exercise}

\begin{exercise}[4]
Find an example of two different metrics on the same set that induce the
same topology.
\end{exercise}

\begin{exercise}[4]
Find an example of two metrics on the same set that induce different
topologies.
\end{exercise}

\begin{exercise}[4]
Research the topologist's sine curve and explain why it is connected but not
path connected.
\end{exercise}

\begin{exercise}[4]
Investigate an example in which sequence convergence does not completely
determine the topology.
\end{exercise}

\begin{exercise}[4]
Research nets and explain why they generalize sequences in arbitrary
topological spaces.
\end{exercise}

\begin{exercise}[4]
Research filters and compare them with nets.
\end{exercise}

\begin{exercise}[4]
Investigate one metrization theorem and describe the conditions under which
a topology is induced by a metric.
\end{exercise}

%%%%%%%%%%%%%%%%%%%%%%%%%%%%%%%%%%%%%%%%%%%%%%%%%%%%%%%%%%%%
\subsection{Conceptual Summary}
\label{subsec:point-set-topology-summary}
%%%%%%%%%%%%%%%%%%%%%%%%%%%%%%%%%%%%%%%%%%%%%%%%%%%%%%%%%%%%

A topological space is a pair

\[
\boxed{
(X,\mathcal{T}),
}
\]

where \(\mathcal{T}\) satisfies the open-set axioms.

The topology determines which subsets are open.

Closed sets are complements of open sets.

For a subset \(A\subseteq X\),

\[
\boxed{
\operatorname{int}(A)
}
\]

is its largest open subset,

\[
\boxed{
\overline{A}
}
\]

is its smallest closed superset, and

\[
\boxed{
\partial A
=
\overline{A}
\setminus
\operatorname{int}(A)
}
\]

is its boundary.

A subset is dense when

\[
\boxed{
\overline{A}=X.
}
\]

Bases and subbases provide smaller collections from which entire topologies
can be generated.

Important topology constructions include

\[
\boxed{
\text{subspace topology},
\quad
\text{product topology},
\quad
\text{quotient topology}.
}
\]

Continuity is characterized by open-set preimages:

\[
\boxed{
f\text{ continuous}
\iff
f^{-1}(U)\text{ open whenever }U\text{ is open}.
}
\]

A homeomorphism is a continuous bijection with continuous inverse.

Compactness is the finite-subcover property:

\[
\boxed{
\text{open cover}
\Longrightarrow
\text{finite subcover}.
}
\]

Connectedness means that the space cannot be separated into two disjoint,
nonempty open pieces.

Path connectedness requires continuous paths between every pair of points.

Hausdorff spaces allow distinct points to be separated by disjoint open
neighborhoods.

Point-set topology therefore provides the fundamental language of

\[
\boxed{
\text{open sets}
+
\text{continuity}
+
\text{compactness}
+
\text{connectedness}
+
\text{separation}.
}
\]

%%%%%%%%%%%%%%%%%%%%%%%%%%%%%%%%%%%%%%%%%%%%%%%%%%%%%%%%%%%%
\subsection{Section Connections}
\label{subsec:point-set-topology-connections}
%%%%%%%%%%%%%%%%%%%%%%%%%%%%%%%%%%%%%%%%%%%%%%%%%%%%%%%%%%%%

\chapterconnection{
Point-Set Topology establishes the language used throughout the remainder of
Topology. Open sets, neighborhoods, continuity, compactness, connectedness,
quotient spaces, and homeomorphisms provide the foundation for Algebraic
Topology. Local Euclidean structure will be used to define manifolds in
Differential Topology and Manifold Theory. Quotient constructions will
reappear in Geometric Topology, knot spaces, surfaces, and identification
spaces. Metric topology connects directly with Analysis, while compactness
and connectedness explain many classical theorems of calculus and real
analysis. The next section develops Algebraic Topology, where algebraic
objects are assigned to topological spaces in order to detect global
structure.
}

%%%%%%%%%%%%%%%%%%%%%%%%%%%%%%%%%%%%%%%%%%%%%%%%%%%%%%%%%%%%
% SECTION SOURCE TRACKING
%%%%%%%%%%%%%%%%%%%%%%%%%%%%%%%%%%%%%%%%%%%%%%%%%%%%%%%%%%%%

%------------------------------------------------------------
% 4.1 Point-Set Topology
%------------------------------------------------------------
%
% Source basis:
%   - Standard introductory general topology and
%     point-set topology.
%
% Used for:
%   - topological spaces;
%   - topology axioms;
%   - open and closed sets;
%   - discrete and indiscrete topologies;
%   - cofinite topology;
%   - metric-induced topologies;
%   - finer and coarser topologies;
%   - neighborhoods;
%   - interior, closure, exterior, and boundary;
%   - limit and isolated points;
%   - dense subsets;
%   - bases and subbases;
%   - local bases;
%   - first and second countability;
%   - subspace topology;
%   - product topology;
%   - quotient topology;
%   - topological continuity;
%   - homeomorphisms;
%   - topological invariants;
%   - sequence convergence;
%   - Hausdorff spaces;
%   - separation axioms;
%   - compactness;
%   - Heine--Borel theorem;
%   - connectedness;
%   - path connectedness;
%   - connected components;
%   - applications to analysis and geometry.
%
% Notes:
%   - Set operations, power sets, functions, relations, and
%     equivalence relations are developed canonically in
%     Foundations.
%   - Metric-space limits and epsilon-based analysis are
%     developed canonically in Analysis.
%   - The full theory of nets, filters, metrization,
%     paracompactness, and advanced general topology may be
%     expanded in later topology material where needed.
%   - Quotient spaces, paths, connectedness, and
%     homeomorphisms prepare for Algebraic Topology.
%   - Manifolds are developed later in the Topology chapter
%     and connected with Differential Geometry.
%
%%%%%%%%%%%%%%%%%%%%%%%%%%%%%%%%%%%%%%%%%%%%%%%%%%%%%%%%%%%%